\par\medskip
\noindent 2.3.17. Aprovechemos la descripción explícita mediante
cociclos del acoplamiento
\[
 \varphi_\xi^n:{}_nP\otimes{}_nQ\longrightarrow G
\]
dada en (2.3.12), para obtener otra expresión del mismo. Supongamos
primero que la biextensión $E$ de clase $\xi$ admite una sección $s$
sobre $P\times Q$.

%%%%%%%%%%  scan idx 261  |  folio 250  %%%%%%%%%%

Como en VII 3.8.13, tomemos para resolución plana de $P$ la que tiene
$L_0=\mathbb Z[P]$ y $L_1$ igual al núcleo del homomorfismo canónico
$\mathbb Z[P]\longrightarrow P$, y definamos $M_\bullet$ de manera
análoga. Existe entonces una única trivialización de biextensión de la
imagen inversa de $E$ sobre $(L_0,M_0)$ cuya restricción a $P\times Q$
está dada por $s$; denotemos esta trivialización también por $s$
(cf. VII 3.8.2). Por VII 3.8.12, la clase de $\xi$ en
\[
 \Ext^1\!\left(
 P\mathbin{\overset{\mathbb L}{\otimes}}Q,G\right)
\]
está representada por el cociclo de componentes $(f,g)$ caracterizadas
por
\[
 s(x_1,y_0)=f(x_1\otimes y_0)e_{E/Q}(y_0),
 \qquad
 s(x_0,y_1)=g(x_0\otimes y_1)e_{E/P}(x_0).
\]
Si $x$ e $y$ son secciones de $\,{}_nP$ y $\,{}_nQ$, levantémoslas a
las secciones correspondientes $x_0=[x]$ e $y_0=[y]$ de $L_0$ y
$M_0$. Aplicando la fórmula (2.3.13), y poniendo $z=s(x,y)$, se
obtiene
\begin{equation}\tag{2.3.18}
 \varphi_\xi^n(x,y)
 ={}^n z/e_{E/Q}(y)-z^n/e_{E/P}(x).
\end{equation}
Aquí el exponente $n$ escrito a la izquierda, respectivamente a la
derecha, de $z$ denota la potencia $n$-ésima para la estructura
multiplicativa izquierda, respectivamente derecha, de $E$. Da, pues,
una sección de $E$ sobre $(0,y)$, respectivamente $(x,0)$; el sentido
del ``cociente'' de un elemento tal por $e_{E/Q}(y)$, respectivamente
$e_{E/P}(x)$, es claro.

La fórmula precedente (2.3.18) conserva sentido siempre que $z$ sea
una sección global de $E$ que levante las secciones $x,y$ de
$\,{}_nP,\,{}_nQ$. Afirmo que (2.3.18) es válida para toda sección de
este tipo, sin suponer que exista una sección de $E$ sobre $P\times Q$.
La sección $(x,y)$ de $P\times Q$ siempre se levanta localmente a una
sección de $E$, puesto que $E\longrightarrow P\times Q$ es un
epimorfismo; por consiguiente, la fórmula calcula
$\varphi_\xi^n(x,y)$ en todos los casos.

%%%%%%%%%%  scan idx 262  |  folio 251  %%%%%%%%%%

Además, si $z$ se reemplaza por otra sección $z'$ de $E$ que levante
$(x,y)$, entonces $z'=z+g$ para una sección $g$ de $G$. Los dos
términos del miembro derecho de (2.3.18) correspondientes a $z'$ son
los correspondientes a $z$, cada uno aumentado en $ng$. Su diferencia
es, por tanto, independiente de $z$, lo que muestra que (2.3.18)
define un homomorfismo bien determinado
\[
 {}_nP\otimes{}_nQ\longrightarrow G.
\]

Queda por demostrar que este homomorfismo es en efecto
$\sisi{\varphi^n}{\varphi_\xi^n}$. Para ello nos reducimos de inmediato
al caso
\[
 P=Q=(\mathbb Z/n\mathbb Z)_T.
\]
Entonces $P\times Q$ es de la forma $I_T$, donde
$I=\mathbb Z/n\mathbb Z\times\mathbb Z/n\mathbb Z$ es un conjunto
finito. En este caso, el torsor $E$ bajo $G_{P\times Q}$, que se
identifica con una familia de torsores $E_i$ bajo $G$, para $i\in I$,
se hace manifiestamente trivial localmente sobre $T$. Quedamos así
reducidos al caso ya tratado en que $E$ posee una sección global sobre
$P\times Q$.
